\notechapter{N-20220315}{Quadratic Parameters, Tangency, and Constrained Extrema}

\section{A scattered algebraic note}

Elementary-looking algebraic problems often hide a parameter, a tangency condition, or an extremum under constraint.  The problems below are organized around those three mechanisms.

\section{A quadratic with parameters}

Let
\[
  f(x)=ax^2+bx+c,\qquad a<b<c,
\]
and suppose the graph passes through \(A=(1,0)\) and \(B=(m,-a)\).  The condition \(f(1)=0\) gives
\[
  a+b+c=0.
\]
The condition \(f(m)=-a\) is equivalent to saying that \(f(x)+a\) has \(x=m\) as a root:
\[
  am^2+bm+c+a=0.
\]
Using \(c=-a-b\), this becomes
\[
  am^2+bm-b=0,\qquad am^2=-b(m-1).
\]
The axis of symmetry is
\[
  x_0=-\frac{b}{2a}.
\]
Thus the parameter \(m\) and the axis \(x_0\) are tied by
\[
  x_0=\frac{m^2}{2m-2}.
\]
Equivalently,
\[
  m^2=(2m-2)x_0,\qquad m^2-2x_0m+2x_0=0.
\]
Solving for \(m\) gives
\[
  m=x_0\pm\sqrt{x_0^2-2x_0}.
\]
Use this relation to examine endpoints and stationary points of \(m\) as a function of \(x_0\).  The extremal parameter occurs at
\[
  x_0=0\quad\text{or}\quad x_0=-\frac12,
\]
and the relevant value is
\[
  m=\frac{\sqrt5-1}{2}.
\]
This is a small example of a recurring trick: convert the algebraic parameter into a geometric parameter, usually the axis, vertex, or tangency point.

\section{A parabola and a hyperbola}

The second problem considers the parabola
\[
  y=x^2-2ax+a^2-4=(x-a)^2-4
\]
and the rectangular hyperbola
\[
  y=\frac{k}{x},\qquad k>0,
\]
in the first quadrant.  The intersection equation is
\[
  x^2-2ax+a^2-4=\frac{k}{x},\qquad x>0,
\]
or
\[
  x^3-2ax^2+(a^2-4)x-k=0.
\]
Questions about ``only one intersection'' are therefore tangency or root-counting questions for this cubic, but the geometry is usually cleaner: one restricts to the right side of the parabola, because the branch \(k/x\) lies in the first quadrant.

\begin{remark}
The hidden principle is that a graph-intersection problem may become either a polynomial root-counting problem or a tangency problem.  The latter is often more geometric and less computational.
\end{remark}

\section{A constrained extremum}

The third problem asks for extrema of
\[
  x+y
\]
under the constraint
\[
  x^2+y^3=2.
\]
Using Lagrange multipliers for
\[
  F(x,y,\lambda)=x+y+\lambda(x^2+y^3-2),
\]
we get
\[
  1+2\lambda x=0,\qquad
  1+3\lambda y^2=0,\qquad
  x^2+y^3=2.
\]
Eliminating \(\lambda\) gives
\[
  2x=3y^2.
\]
Hence
\[
  x=\frac32y^2.
\]
Substitution into the constraint gives
\[
  \frac94y^4+y^3=2,
\]
so the stationary points come from the quartic
\[
  9y^4+4y^3-8=0.
\]
Numerical solution gives approximately
\[
  A=(1.15195,2.02829),\qquad B=(1.82947,0.72509),
\]
for the function
\[
  h(x)=x+(2-x^2)^{1/3}.
\]
These are the local maximum and local minimum along the real branch.

\begin{figure}[H]
\centering
\begin{tikzpicture}[scale=1.2,>=Stealth,every node/.style={fill=white,inner sep=1.2pt}]
  \draw[->] (-1.5,0)--(4.3,0) node[right] {$x$};
  \draw[->] (0,-1.25)--(0,3.2) node[above] {$y$};
  \draw[smooth,thick] plot coordinates {(-1.2,-0.45) (-0.5,0.55) (0,1.26) (0.6,1.78) (1.15,2.03) (1.45,1.55) (1.83,0.73) (2.5,0.95) (4.0,1.65)};
  \fill (1.15,2.03) circle (0.04);
  \node[above left=2pt] at (1.15,2.03) {$A\approx(1.15,2.03)$};
  \fill (1.83,0.73) circle (0.04);
  \node[below right=2pt] at (1.83,0.73) {$B\approx(1.83,0.73)$};
  \node[below] at (2.25,-0.55) {$y=x+(2-x^2)^{1/3}$};
\end{tikzpicture}
\caption{Reconstruction of the graph used for the constrained extremum problem.}
\end{figure}


\section{Tangency as a double-root condition}

A parameter problem involving a quadratic often asks when a line is tangent to a parabola or when an equation has exactly one solution.  Algebraically, this is a discriminant condition:
\[
  \Delta=b^2-4ac=0.
\]
Geometrically, the line touches the parabola at one point.  The same idea appears in constrained extrema: a level curve just touches a constraint curve at an optimum.

\section{Completing the square}

For a quadratic
\[
  ax^2+bx+c,\qquad a\ne0,
\]
completing the square gives
\[
  a\left(x+\frac b{2a}\right)^2+c-\frac{b^2}{4a}.
\]
This form reveals the vertex, minimum or maximum value, and symmetry axis.  It is usually more informative than the expanded form.

\section{A wider frame}

The three examples are connected by the same idea: reduce a complicated statement to a structural condition.  The axis of a quadratic, the tangency of two graphs, and the Lagrange multiplier equations are all ways of saying that an extremum occurs when a first-order freedom disappears.

\section{Tangency as multiplicity}

A line is tangent to a curve when the intersection equation has a repeated root.  If a family is written as
\[
  F(x,t)=0,
\]
then a tangency or transition value often satisfies
\[
  F(x,t)=0,\qquad \frac{\partial F}{\partial x}(x,t)=0.
\]
For a quadratic, this is the discriminant condition \(\Delta=0\).  The same idea persists for higher-degree equations: a repeated root is detected by a common factor of \(F\) and \(F_x\).

\section{Constrained extrema and normal directions}

For a constraint \(g(x,y)=0\), the Lagrange multiplier condition
\[
  \nabla f=\lambda \nabla g
\]
says that at an interior constrained extremum, the gradient of \(f\) has no component tangent to the constraint curve.  Equivalently, the differential \(df\) vanishes on the tangent space of the constraint.

This is the coordinate-free meaning of the method: extrema occur where the objective's first-order change is normal to every allowed first-order motion.

\section{Degeneracy and second-order data}

After first derivatives vanish, the next question is second-order behavior.  For an unconstrained two-variable function, the Hessian quadratic form determines the local model:
\[
  f(a+h)=f(a)+\frac12 h^T H_f(a)h+o(\norm h^2).
\]
For constrained problems, the Hessian must be restricted to tangent directions.  Degenerate cases occur when this quadratic form has zero directions, forcing higher-order analysis.

\section{Multiplicity, normality, and local geometry}

Quadratic parameter problems, tangency, and constrained extrema are one theme: detect where the number or type of solutions changes.  Algebra sees this as multiplicity or discriminant zero; geometry sees it as tangency; optimization sees it as vanishing first variation plus second-order testing.
